% Collision-energy audit for CODEX_SPEC_energy.md.
% This is a fragment intended to be read with proof.tex, nv_theorem.tex,
% and meso_result.tex.  The finite computations are in energy_verify.py.

\subsection{Collision energy is an off-boundary factorial moment}

Retain
\[
 \mathscr D_H=\{(d,r)\in\mathbb Z_{\geq2}^2:d+r\leq H\},
 \qquad
 \mathcal E_p(H)=\sum_{(d,r)\in\mathscr D_H}m_{d,r}(p),
\]
and assume throughout that $p\geq7$ and $H\leq\sqrt p$; the omitted
small cases are empty or immediate.  For $x\in\mathbb F_p$ put
\[
 k_x=k_x(H)=\#\{2\leq h\leq H:N_h(x)=0\},
 \qquad R_p(H)=\sum_{x\in\mathbb F_p}k_x,
\]
and put
\[
 \partial_H=\{-2,-3,\ldots,-H\}\subseteq\mathbb F_p.
\]
Split the energy as
\[
 \mathcal E_p(H)=\mathcal E_p^\circ(H)+\mathcal E_p^\partial(H)
\]
according as its root witness lies outside or inside $\partial_H$.

\begin{proposition}[Exact off-boundary column identity]
\label{prop:energy-column-exact}
One has
\begin{equation}\label{eq:energy-column-exact}
 \boxed{\mathcal E_p^\circ(H)
 =\sum_{x\notin\partial_H}\binom{k_x}{2}.}
\end{equation}
At the boundary only the following direction is universally valid:
\begin{equation}\label{eq:energy-boundary-sandwich}
 0\leq\mathcal E_p^\partial(H)
 \leq\sum_{x\in\partial_H}\binom{k_x}{2}.
\end{equation}
Consequently
\begin{equation}\label{eq:energy-global-sandwich}
 \sum_{x\notin\partial_H}\binom{k_x}{2}
 \leq\mathcal E_p(H)
 \leq\sum_{x\in\mathbb F_p}\binom{k_x}{2}.
\end{equation}
\end{proposition}

\begin{proof}
The continuant addition formula used in the proof of
Proposition~\ref{prop:bezout} gives
\[
 N_{d+r}(x)
 =N_d(x)N_{r+1}(x+d-1)
  -(x+d)^6N_{d-1}(x)N_r(x+d).
\]
Thus an energy witness $(x,d,r)$ gives the two zero levels $d$ and
$d+r$ in the column $x$.  This map is injective, since its lower level
and the difference recover $(d,r)$.  It proves the upper bound in
\eqref{eq:energy-boundary-sandwich} and the upper bound in
\eqref{eq:energy-global-sandwich}.

Conversely, suppose that $x\notin\partial_H$ and that $2\leq d<e\leq H$
are two zero levels.  Then $S_d(x)$ is a unit, so the saturated B\'ezout
identity gives
\[
 N_{e-d}(x+d)=0.
\]
The difference $e-d$ cannot be one: in that case the right-hand side of
\eqref{eq:bezout} would be $S_d(x)^6N_1=S_d(x)^6\ne0$.  Hence
$(x,d,e-d)$ is an energy witness.  This is inverse to the preceding map
and proves~\eqref{eq:energy-column-exact}.
\end{proof}

The equality is stronger than~\eqref{eq:meso-column-pairs}, but it also
shows why the boundary exclusion cannot be repaired by changing an
inequality sign.  At $p=73$, $H=8$, the polluted column $x=-3$ has six
zero levels and hence $15$ column pairs, while only two boundary energy
witnesses occur in the entire computation.

\begin{corollary}[What the present column estimates prove]
\label{cor:energy-eight-thirds}
Uniformly for $H\leq\sqrt p$,
\begin{equation}\label{eq:energy-eight-thirds}
 \mathcal E_p(H)\ll H^{8/3}.
\end{equation}
\end{corollary}

\begin{proof}
The row-degree bound gives
\begin{equation}\label{eq:energy-R-trivial}
 R_p(H)\leq\sum_{h=2}^H3(h-1)=\frac32H(H-1).
\end{equation}
Define the window-polluted set
\[
 \mathcal P_H=
 \{-m:2\leq m\leq H,\ p\mid b_{m-1}\}.
\]
The pollution characterization following Lemma~\ref{lem:restart} and
Corollary~\ref{cor:subinterval} give
\[
 P_H:=|\mathcal P_H|\ll H^{2/3}.
\]
Every column outside $\mathcal P_H$ has no adjacent zero levels in
$[2,H]$.  (A globally polluted column whose cut occurs after $H$ has no
adjacent pair in this window.)  Proposition~\ref{prop:column} therefore
gives
\[
 M_H:=\max_{x\notin\mathcal P_H}k_x\ll H^{2/3}.
\]
By Proposition~\ref{prop:energy-column-exact},
\begin{align*}
 \mathcal E_p(H)
 &\leq\sum_x\binom{k_x}{2}\\
 &\leq\frac{M_H-1}{2}
       \sum_{x\notin\mathcal P_H}k_x
       +P_H\binom{H-1}{2}\\
 &\ll H^{2/3}H^2+H^{2/3}H^2.
\end{align*}
This is~\eqref{eq:energy-eight-thirds}.
\end{proof}

This saves $H^{1/3}$ over the direct $O(H^3)$ degree bound, but its
exponent is still larger than two and hence has no consequence for the
record exponent of $Z(p)$.

\subsection{The proposed fixed point is not a contraction}

For an arbitrary threshold $2\leq K\leq H$, define
\[
 R_p^\circ(H)=\sum_{x\notin\partial_H}k_x,
 \qquad
 L_{p,K}(H)=
 \sum_{\substack{x\notin\partial_H\\0<k_x<K}}k_x,
\]
and let
\[
 M_p^\circ(H)=\max_{x\notin\partial_H}k_x.
\]

\begin{proposition}[Exact fixed-point audit]
\label{prop:energy-fixed-point-audit}
For every such $K$,
\begin{equation}\label{eq:energy-amplification-K}
 R_p(H)\leq R_p^\partial(H)+L_{p,K}(H)
 +\frac{2}{K-1}\mathcal E_p^\circ(H),
\end{equation}
where $R_p^\partial(H)=\sum_{x\in\partial_H}k_x$.  On the other hand,
\begin{equation}\label{eq:energy-upper-from-R}
 \mathcal E_p^\circ(H)
 \leq\frac{M_p^\circ(H)-1}{2}R_p^\circ(H).
\end{equation}
Combining the two estimates never gives an absorptive fixed point.
\end{proposition}

\begin{proof}
For $k\geq K$ one has
\[
 k\leq\frac{2}{K-1}\binom{k}{2}.
\]
Summing this over the high off-boundary columns and using
\eqref{eq:energy-column-exact} proves~\eqref{eq:energy-amplification-K}.
The elementary inequality
$\binom{k}{2}\leq(M_p^\circ-1)k/2$ proves
\eqref{eq:energy-upper-from-R}.

Substitution gives
\begin{equation}\label{eq:energy-feedback}
 \mathcal E_p^\circ
 \leq\frac{M_p^\circ-1}{2}L_{p,K}
 +\frac{M_p^\circ-1}{K-1}\mathcal E_p^\circ.
\end{equation}
If $K\leq M_p^\circ$, the feedback coefficient is at least one.  If
$K>M_p^\circ$, every nonzero off-boundary column is already included in
$L_{p,K}$, so $L_{p,K}=R_p^\circ$ and~\eqref{eq:energy-feedback} is the
original tautology.  These two cases exhaust all thresholds.
\end{proof}

In particular, the proposed $K=\lceil\sqrt H\rceil$ does not interact
contractively with the available $M_p^\circ\ll H^{2/3}$ estimate.  The
amplification inequality is a lower bound on the factorial moment carried
by high fibers; it is not a second upper bound for that moment.  Singleton
columns can make $R_p(H)$ large while contributing nothing to energy, and
even an independent $R_p(H)\ll H$ estimate together with only
$k_x\ll H^{2/3}$ permits the abstract scale
$\mathcal E_p(H)\asymp H^{5/3}$.

There is also a notation obstruction in the suggested partition.  In
\texttt{meso\_result.tex}, $W_p(H)$ denotes simple affine support and
$\mathcal W_p(H)$ denotes degree-weighted raw resultant support.  A raw
resultant-supported pair need have no $\mathbb F_p$ witness at all, and
therefore cannot be assigned to a witness column.  Even after restricting
to affine support, the artificial weight $\min(d,r)-1$ is not bounded by
the number of actual witnesses.  Thus an inequality of the form
``$\mathcal W_p(H)\leq\sum_x\binom{k_x}{2}$'' is unavailable; the valid
column statement is Proposition~\ref{prop:energy-column-exact} for the
actual energy.

\subsection{The exact triple-counting exponent}

\begin{proposition}[Optimization from an energy exponent]
\label{prop:energy-beta-optimization}
Suppose, uniformly in a range of $H$, that
\[
 \mathcal E_p(H)\ll H^\beta.
\]
Then block triple-counting gives
\begin{equation}\label{eq:energy-beta-endpoint}
 Z(p)\ll\frac pH+p^{2/3}H^{(\beta-2)/3}.
\end{equation}
If the estimate is uniform for every $H\leq\sqrt p$, the optimized
triple-counting exponent is
\begin{equation}\label{eq:energy-beta-piecewise}
 \theta_{\rm triple}(\beta)=
 \begin{cases}
  1/2,&\beta\leq1,\\
  (\beta+2)/6,&1\leq\beta\leq2,\\
  \beta/(\beta+1),&\beta\geq2.
 \end{cases}
\end{equation}
After combining with the existing $Z(p)\ll p^{2/3}$ estimate, every
$1<\beta<2$ gives
\[
 Z(p)\ll p^{(\beta+2)/6},
\]
and in particular $\beta=3/2$ gives $Z(p)\ll p^{7/12}$.
\end{proposition}

\begin{proof}
Partition the index interval into $B\ll p/H$ nonwrapping blocks and let
$z_i$ be their zero counts.  As in the proof of
\eqref{eq:meso-triple-endpoint},
\[
 \sum_i\binom{z_i}{3}\leq\mathcal E_p(H).
\]
The elementary bound
\[
 z\leq2+\bigl(6\tbinom z3\bigr)^{1/3}
\]
and H\"older's inequality give
\[
 Z(p)\ll B+B^{2/3}\mathcal E_p(H)^{1/3},
\]
which is~\eqref{eq:energy-beta-endpoint}.

Write $H=p^\alpha$, $0<\alpha\leq1/2$.  The two exponents are
\[
 1-\alpha,
 \qquad
 \frac23+\frac{\beta-2}{3}\alpha.
\]
When $\beta<2$, both decrease with $\alpha$, so one takes
$\alpha=1/2$.  This gives
$\max\{1/2,(\beta+2)/6\}$.  At $\beta=2$ the second exponent is $2/3$.
For $\beta>2$, balancing the two terms at
$\alpha=1/(\beta+1)$ gives $\beta/(\beta+1)$.  This proves
\eqref{eq:energy-beta-piecewise} and the stated comparison.
\end{proof}

Thus the power $(\beta-1)/3$ displayed in the specification is not the
output of block triple-counting: the factor $B^{2/3}$ contributes
$H^{-2/3}$, not $H^{-1/3}$.  Equivalently, the displayed formula is
missing one further factor $H^{-1/3}$.  More generally, if
$\beta<2$ is available only through $H\leq p^\eta$, then taking the
largest permitted $H$ gives
\begin{equation}\label{eq:energy-range-beta}
 \theta(\beta,\eta)=
 \max\left\{1-\eta,
 \frac23+\frac{\eta(\beta-2)}3\right\}.
\end{equation}
It beats $2/3$ exactly when $\eta>1/3$ and $\beta<2$.  No interior
balancing is required in this decreasing regime.

\subsection{Structured subfamilies}

To avoid the clash between simple affine support and weighted resultant
support, write, for $\mathscr S\subseteq\mathscr D_H$,
\[
 \mathcal W_p^{\rm res}(\mathscr S)
 =\sum_{(d,r)\in\mathscr S}
  (\min(d,r)-1)\mathbf1_{p\mid\mathcal R_{d,r}},
 \qquad
 \mathcal E_p(\mathscr S)
 =\sum_{(d,r)\in\mathscr S}m_{d,r}(p).
\]

\begin{proposition}[Shallow strips and diagonal bands]
\label{prop:energy-structured-combinatorial}
Let $2\leq K\leq H/2$ and $k=\lfloor K\rfloor$.  Then
\begin{align}
 \mathcal W_p^{\rm res}
 \{(d,r):\min(d,r)\leq K\}
 &\leq B_H(k),\label{eq:energy-shallow-W}\\
 \mathcal E_p
 \{(d,r):\min(d,r)\leq K\}
 &\leq3B_H(k),\label{eq:energy-shallow-E}
\end{align}
where
\begin{equation}\label{eq:energy-BHK}
 B_H(k)=\sum_{s=2}^k(s-1)(2H-4s+1)
 \leq Hk(k-1).
\end{equation}
In particular, both quantities are $O(H^{3/2})$ in the deterministic
range
\[
 \min(d,r)\leq H^{1/4}.
\]
The larger suggested strip $\min(d,r)\leq\sqrt H$ has only the bound
$O(H^2)$ after the necessary degree weight is included.

For a fixed $C\geq0$ and $D\leq H/2$,
\begin{equation}\label{eq:energy-diagonal-band-partial}
 \mathcal W_p^{\rm res}
 \{(d,r):|d-r|\leq C,\ \min(d,r)\leq D\}
 \ll(C+1)D^2,
\end{equation}
and the same bound, up to a factor three, holds for the energy.  Hence
the part $\min(d,r)\leq H^{3/4}$ of a fixed-width diagonal band is
$O_C(H^{3/2})$, whereas the full band has only the present bound
$O_C(H^2)$.
\end{proposition}

\begin{proof}
For a fixed value $s=\min(d,r)$ there are exactly
$2H-4s+1$ ordered pairs in $\mathscr D_H$.  Summing the weight $s-1$
gives~\eqref{eq:energy-BHK}.  Since the reductions of $N_h$ are nonzero
in the present range,
\[
 m_{d,r}(p)\leq3(\min(d,r)-1),
\]
which proves~\eqref{eq:energy-shallow-E}.

In the band $|d-r|\leq C$, each possible minimum occurs in at most
$2C+1$ ordered pairs.  Summing its degree weight gives
\eqref{eq:energy-diagonal-band-partial}.
\end{proof}

The thresholds are not artifacts of a loose summation: exactly
\[
 B_H(k)=Hk(k-1)+O(k^3).
\]
Thus if $k\to\infty$ and $k=o(H)$, then
$B_H(k)=(1+o(1))Hk^2$.  The unweighted support of the
$\sqrt H$ strip is $O(H^{3/2})$, but this says nothing about its energy
multiplicities.  Even after discarding shifted columns which are
polluted, Proposition~\ref{prop:column} gives only
\begin{equation}\label{eq:energy-unpolluted-shallow}
 \sum_{\substack{2\leq d\leq K,\ 2\leq r,\ d+r\leq H\\
                  x\in X_{d,r}(p),\ x+d\ {\mathrm{unpolluted}}}}1
 \ll K^2H^{2/3};
\end{equation}
The symmetric orientation obeys the same estimate.  For $K=\sqrt H$
this is $H^{5/3}$, still above the target.

The ambient degree weight satisfies
\begin{equation}\label{eq:energy-total-ambient-weight}
 \sum_{(d,r)\in\mathscr D_H}(\min(d,r)-1)
 =\frac{H^3}{12}+O(H^2).
\end{equation}
Removing the whole $\sqrt H$ strip and any fixed-width diagonal band
changes~\eqref{eq:energy-total-ambient-weight} by only $O_C(H^2)$.
Thus the ambient dominant region is the balanced off-diagonal bulk
\[
 \min(d,r)>\sqrt H,\qquad |d-r|>C.
\]
This does not mean that the actual worst contribution has been localized:
the shallow annulus $H^{1/4}<\min(d,r)\leq\sqrt H$ and the high part of
the diagonal band are themselves unresolved at the $H^{3/2}$ scale.

\subsubsection{Infinity and center lattices}

\begin{proposition}[Exact weight of the infinity lattice]
\label{prop:energy-infinity-weight}
Let $\rho=\rho_p$ and $M=\lfloor H/\rho\rfloor$.  The automatic
projective-infinity pairs $\rho\mid d,r$ have raw degree weight
\begin{equation}\label{eq:energy-infinity-weight}
 \mathcal I_p(H)
 =\rho\sum_{s=2}^{M}\left\lfloor\frac{s^2}{4}\right\rfloor
  -\binom M2
 =\frac{H^3}{12\rho^2}+O\left(\frac{H^2}{\rho}\right).
\end{equation}
They contribute no automatic affine witness and must be saturated before
using resultant support as a collision proxy.
\end{proposition}

\begin{proof}
Write $d=\rho a$, $r=\rho b$.  For $s=a+b$,
\[
 \sum_{a=1}^{s-1}\min(a,s-a)=\left\lfloor\frac{s^2}{4}\right\rfloor.
\]
There are $\binom M2$ ordered pairs, and subtracting the $1$ in each
weight gives~\eqref{eq:energy-infinity-weight}.

Let $\overline N_d$ be the reduction of $N_d$ modulo $p$, let
$n_d=3(d-1)$ be its nominal degree, and define its nominal
homogenization by
\[
 \overline F_d(X,Z)=Z^{n_d}\overline N_d(X/Z).
\]
When $\rho\mid d$ and $d<p$, the centered-coefficient calculation after
\eqref{eq:meso-rank-lattice} shows that the nominal homogenization has
the exact form
\[
 \overline F_d(X,Z)=Z^2F_d^{\rm sat}(X,Z),
 \qquad F_d^{\rm sat}(1,0)\ne0.
\]
Modulo $p$, dividing $Z^2$ from each nominal form removes the automatic
infinity intersection exactly.  One must not delete the whole pair: its
saturated forms may still have a separate affine common root.
\end{proof}

For example, the already recorded values $p=665857$, $\rho_p=8$, and
$H=816$ give $\mathcal I_p(H)=712657$, whereas
$H^{3/2}=23309.6\ldots$.  This is why the raw weighted resultant support
is not the statement to which a witness-column proof can apply.

The center factors have the opposite status: they are genuine affine
witnesses and cannot be saturated away.

\begin{proposition}[Present bound for distinguished center witnesses]
\label{prop:energy-center-witnesses}
For $(d,r)\in\mathscr D_H$, define
\[
 c_{d\to r}(p)=
 \begin{cases}
  \mathbf1_{p\mid T_r^{(d)}},&d\ \text{even},\\
  0,&d\ \text{odd},
 \end{cases}
 \qquad
 c_{d,r}(p)=c_{d\to r}(p)+c_{r\to d}(p).
\]
Then $c_{d,r}(p)$ is the number of distinguished central
$\mathbb F_p$-witnesses supplied by the two even blocks, and
\begin{equation}\label{eq:energy-center-count}
 \sum_{(d,r)\in\mathscr D_H}c_{d,r}(p)\ll H^{5/3}.
\end{equation}
Their degree-weighted support is consequently at most $O(H^{8/3})$ with
the present inputs.
\end{proposition}

\begin{proof}
If $d$ is even, $x=-(d+1)/2$ is a root of $N_d$, and
\[
 N_r(x+d)=N_r((d-1)/2)=0
 \quad\Longleftrightarrow\quad p\mid T_r^{(d)},
\]
because $p$ is odd.  The symmetric root is
$x=-d-(r+1)/2$.  When both occur, their difference is $(d+r)/2\ne0$
modulo $p$, so they are distinct.

For fixed even $d$, the column $y=(d-1)/2$ is nonsingular throughout
$2\leq r\leq H-d$: indeed
\[
 0<2(y+r)=d-1+2r<2H<p.
\]
Proposition~\ref{prop:column} bounds its zero levels by $O(H^{2/3})$.
Summing over the $O(H)$ even values of $d$ and adding the symmetric
orientation proves~\eqref{eq:energy-center-count}.  Multiplication by
the largest possible degree weight gives $O(H^{8/3})$.
\end{proof}

Thus the center structure explains repeated genuine factors, but the
proved column exponent leaves a gap of $H^{1/6}$ even for their
unweighted witness count.

\subsubsection{What the diagonal square does and does not buy}

\begin{proposition}[Diagonal midpoint and split off-center roots]
\label{prop:energy-diagonal-decomposition}
Put
\[
 M_p(H)=\#\{2\leq d\leq H/2:p\mid D_d\}.
\]
Then $M_p(H)\ll H^{2/3}$.  If
\[
 q_d(p)=\#\{u\in(\mathbb F_p^\times)^2:
                   E_d(u)=O_d(u)=0\},
\]
then the diagonal energy has the exact decomposition
\begin{equation}\label{eq:energy-diagonal-decomposition}
 \mathcal E_p^\Delta(H):=\sum_{2d\leq H}m_{d,d}(p)
 =\sum_{2\leq d\leq H/2}
   \left(\mathbf1_{p\mid D_d}+2q_d(p)\right)
 =O(H^{2/3})+2\sum_{2\leq d\leq H/2}q_d(p).
\end{equation}
The corresponding degree-weighted support caused by the $D_d$ factors
is $O(H^{5/3})$.  The unresolved diagonal contribution is the sum of
split off-center counts $q_d(p)$.  Each such root forces $p\mid Q_d$,
but the converse need not hold.
\end{proposition}

\begin{proof}
Since $p$ is odd,
\[
 p\mid D_d\quad\Longleftrightarrow\quad N_d(-1/2)=0.
\]
The column $x=-1/2$ is nonsingular, so
Proposition~\ref{prop:column} gives $M_p(H)\ll H^{2/3}$.  Formula
\eqref{eq:meso-diagonal-root-count} is exactly
\eqref{eq:energy-diagonal-decomposition}.  Weighting the
$O(H^{2/3})$ midpoint indices by at most $H$ gives the last support
bound.
\end{proof}

In the factorization
$\mathcal R_{d,d}=\ell_d|D_d|Q_d^2$, the explicit factor $\ell_d$
comes from the leading coefficient and $D_d$ detects the midpoint
$T=0$; $Q_d$ is the remaining parity resultant.  Its square reflects
the formal parity pairing $T,-T$, which may coincide or degenerate, not
two independent supported indices.
The three factors need not be coprime, and $p\mid Q_d$ does not imply a
split affine root: it may instead come from a nonsplit common $u$, or,
when $p\mid\ell_d$, from a degree degeneration at $U$-infinity.  For
example, $p=577$, $d=4$ has $p\mid\ell_4Q_4$ but no affine common root.

\begin{proposition}[The valuation gain is only a constant in averaging]
\label{prop:energy-diagonal-valuation-average}
Let $H\leq\sqrt X$ and define
\[
 \mathscr G_p(H)=\{2\leq d\leq H/2:
 p\mid\mathcal R_{d,d},\ p\nmid\ell_dD_d\}.
\]
Then
\begin{align}
 \sum_{X<p\leq2X}\#\mathscr G_p(H)
 &\ll\frac{H^3\log H}{\log X},
 \label{eq:energy-diagonal-average-count}\\
 \sum_{X<p\leq2X}\sum_{d\in\mathscr G_p(H)}(d-1)
 &\ll\frac{H^4\log H}{\log X}.
 \label{eq:energy-diagonal-average-weight}
\end{align}
\end{proposition}

\begin{proof}
The diagonal valuation law gives $p^2\mid\mathcal R_{d,d}$ for every
indicated pair.  Therefore
\[
 2\log X\sum_{X<p\leq2X}\#\mathscr G_p(H)
 \leq\sum_{2\leq d\leq H/2}\log\mathcal R_{d,d}.
\]
The Hadamard estimate
$\log\mathcal R_{d,d}\ll d^2\log(2d)$ proves
\eqref{eq:energy-diagonal-average-count}.  Applying the same argument to
$\prod_{2\leq d\leq H/2}\mathcal R_{d,d}^{d-1}$ proves
\eqref{eq:energy-diagonal-average-weight}.
\end{proof}

Without the square valuation, the same argument merely omits the factor
$2$ on the left.  Thus it improves a prime-averaged constant, not an
$H$-exponent.  At the mesoscopic endpoint $H\asymp\sqrt X$, both
resulting estimates are weaker than the trivial $O(H)$ count and
$O(H^2)$ degree weight.  For much shorter ranges they can exclude an
initial polylogarithmic segment, but they give no uniform $H$-power gain
in the target range.  Hence the diagonal square law does not prove the
requested diagonal-band case.

\subsection{The exact remaining lemma and stall report}

For $(d,r)\in\mathscr D_H$, define the split affine gcd degree
\[
 g_{d,r}^{\rm sp}(p)=
 \deg\gcd_{\mathbb F_p[X]}
 \bigl(\overline N_d(X),\overline N_r(X+d),X^p-X\bigr).
\]
The polynomial $X^p-X$ is squarefree, and $N_d,N_r$ are nonzero in the
present range, so
\begin{equation}\label{eq:energy-split-gcd-equals-m}
 g_{d,r}^{\rm sp}(p)=m_{d,r}(p).
\end{equation}
If
\[
 A_p(H;j)=\#\{(d,r)\in\mathscr D_H:g_{d,r}^{\rm sp}(p)\geq j\},
\]
then the layer-cake identity is
\begin{equation}\label{eq:energy-layer-cake}
 \boxed{\mathcal E_p(H)=\sum_{j\geq1}A_p(H;j).}
\end{equation}

Accordingly, the minimal missing arithmetic statement can be isolated as
the following split affine gcd-tail lemma:
\begin{equation}\label{eq:energy-missing-gcd-tail}
 \boxed{
 \sum_{j\geq1}
 \#\left\{(d,r)\in\mathscr D_H:
 \deg\gcd\bigl(\overline N_d,\overline N_r(\,\cdot+d),X^p-X\bigr)
 \geq j\right\}
 \ll H^{3/2}.}
\end{equation}
It is exactly equivalent to the energy target, automatically discards
the infinity lattice and nonsplit common factors, and charges the actual
multiplicity instead of the worst possible degree.  A slightly stronger
purely algebraic sufficient statement replaces the three-polynomial gcd
by the degree of the radical of the affine gcd
$\gcd(\overline N_d,\overline N_r(\,\cdot+d))$.

The degree-weighted raw resultant estimate is a convenient sufficient
condition, but it is not minimal: it both inserts the worst-case degree
$3(\min(d,r)-1)$ and counts projective and nonsplit support.  Likewise,
the fixed-point inequalities contain no mechanism capable of proving
\eqref{eq:energy-missing-gcd-tail}; a new arithmetic incidence or
split-gcd tail estimate is required.

\paragraph{Finite verification.}
The script \texttt{energy\_verify.py} evaluates the continuant recurrence
pointwise for all $669$ primes $p\leq5000$, always with
$H=\lfloor\sqrt p\rfloor$.  It computes $R_p(H)$, total and boundary
energy, both factorial column moments, $R_p^\partial(H)$, low-fiber mass,
and polluted-column counts.  It verifies
\eqref{eq:energy-column-exact}, \eqref{eq:energy-global-sandwich}, and
the integer form of~\eqref{eq:energy-amplification-K} for
$K=\lceil\sqrt H\rceil$ whenever $H\geq2$ (the $667$ primes $p\geq5$).
The column identities are checked for all $669$ primes.  Selected records
are:
\[
\begin{array}{c|r|r|r|r|r|r|r|r}
p&H&R&\mathcal E&\mathcal E^\circ&
 \sum_x\binom{k_x}{2}&R^\partial&P_H&
 \sum_{x\in\mathcal P_H}k_x\\ \hline
73&8&19&2&0&16&11&1&6\\
131&11&17&6&6&6&0&0&0\\
653&25&38&16&16&16&0&0&0\\
3331&57&187&0&0&1225&98&1&50\\
4283&65&203&4&2&1084&93&1&47
\end{array}
\]
The canonical comma-separated $669$-record digest printed by the script
is recorded in the verifier itself.  Across the scan,
\[
 \max\mathcal E_p(H)=16\quad(p=653),
 \qquad
 \max\frac{\mathcal E_p(H)}{H^{3/2}}
 =\frac6{11^{3/2}}\quad(p=131),
\]
and the largest unpolluted column has four returns.  These data are only
tests of the deterministic identities, not evidence used in their proofs.
In fact every off-boundary column in the $667$ records with $K\geq2$ has
$k_x<K$, so $R_p(H)=R_p^\partial(H)+L_{p,K}(H)$ there.  The tested
amplification inequality is therefore completely vacuous on every record
in its stated domain, exactly as the fixed-point audit predicts.

The dynamic-height scan also corrects one reconnaissance datum in the
existing text.  Since
\[
 b_7=584307365=3331\cdot175415,
\]
the column $x=-8$ at $p=3331$ vanishes at every level
$8\leq h\leq57$ and has $k_{-8}=50$.  Thus the statement that the
largest polluted count through $p=5000$ is $19$ at $p=541$ is false for
$H=\lfloor\sqrt p\rfloor$; $19$ is the value at that particular prime.

\paragraph{STALL REPORT.}
The target $\mathcal E_p(H)\ll H^{3/2}$ is not proved.  The rigorous
outputs are:
\begin{enumerate}
 \item the exact off-boundary identity
       \eqref{eq:energy-column-exact} and the sharp boundary directions;
 \item the proof that the proposed $\mathcal E\leftrightarrow R$ fixed
       point cannot contract, together with the unconditional
       $\mathcal E_p(H)\ll H^{8/3}$ consequence of current inputs;
 \item the corrected optimization
       $Z(p)\ll p/H+p^{2/3}H^{(\beta-2)/3}$, under which every
       $\beta<2$ available past $H=p^{1/3}$ improves $2/3$;
 \item $H^{3/2}$ bounds for the deterministic shallow strip
       $\min(d,r)\leq H^{1/4}$ and for the low part
       $\min(d,r)\leq H^{3/4}$ of a fixed-width diagonal band;
 \item exact removal of the infinity lattice, a bound for genuine center
       witnesses, and the isolation of the actual split off-center counts
       $q_d(p)$, each of which forces $p\mid Q_d$;
 \item the exact missing split affine gcd-tail statement
       \eqref{eq:energy-missing-gcd-tail}.
\end{enumerate}
